\honest{Rationality and the English Language}{Eliezer Yudkowsky}{2007}

A commenter said my writing reminded them of George Orwell's ``Politics and the English Language.'' I was honored, and I had already planned today's topic.

Orwell is ``mandatory reading for rationalists as well as authors.'' If you want to imprison people for years without trial, you do not say ``I'm going to imprison Mr. Jennings for years without trial.'' You say, ``Unreliable elements were subjected to an alternative justice process.'' \nb{The example is Orwell's, adapted. His plain version is ``People are imprisoned for years without trial''; his euphemism is ``elimination of unreliable elements.''}

``Writers are told to avoid usage of the passive voice.'' I wrote that sentence in the passive to show the flaw: it does not say who tells them. \nb{One of them is Orwell, whose fourth rule is ``Never use the passive where you can use the active.''} ``Passive voice removes the actor, leaving only the acted-upon.'' With enough static noun phrases, nothing unpleasant ever happens. \nb{Only a passive without a ``by'' phrase does that, and an active sentence can hide the actor too. Orwell's plain version above is itself a passive with no actor; in his example what hides the facts is the noun phrase.} Scientists write ``The subjects were administered Progenitorivox'' because it sounds more authoritative than saying who handed a bottle to whom. \nb{Some scientific bodies asked for the passive. The Acoustical Society of America told authors of meeting abstracts to write ``It was noted'' instead of ``We noted.''}

Someone will object that a reader can work out that the scientist is there. That is why rationalists need Orwell. Nonfiction conveys knowledge; fiction conveys experience. Medicine can predict what happens to a person in a vacuum; fiction can make you live through it.

Charitable readers look for a sensible meaning in a misleading phrase. Authors ``are trained not to give themselves the benefit of the doubt.'' \nb{A commenter asked, ``Trained by whom?'' Yudkowsky's next comment was ``Yes, well, I never said I was a True Writing Master either.''} ``Whatever the audience thinks you said is what you said.'' A writer knows readers will not stop to think; they get a stream of first impressions. Then I apply this rule for writers to other people's statements: judge them by their first impression, not by meanings found through analysis. A novelist would see that the Progenitorivox sentence gives ``the feeling of being told something reliable'' and nothing more, where a concrete version would show the postdoc holding out the bottle and the student's nervous grin. I do not want journal articles written like novels; I want rationalists aware of what words make a reader experience. I report my impression as the only experience any reader has.

A rationalist must become aware of the ``experiential impact of phrases.'' ``Meaning does not excuse impact!'' No interpretation of an applause light such as ``AI should be developed through democratic processes'' excuses its effect. \nb{The example comes from ``Applause Lights,'' posted the day before. Its speaker disputed that post's account in the comments there.}

I close with two passages from Orwell. In the first, a tired speaker repeating stock phrases such as ``BESTIAL, ATROCITIES, IRON HEEL'' seems less a person than a dummy; the noises come out, but the brain is not involved. In the second, Orwell asks us to ``let the meaning choose the word.'' When we think of a concrete object we think without words and then hunt for words that fit. When we think of something abstract, the existing phrases rush in and blur or change the meaning, so it is better to get the meaning clear through pictures and sensations first. \nb{Orwell's next sentence, which the post leaves out, makes the post's point about the audience: one should ``decide what impressions one's words are likely to make on another person.''} ``Charles Sanders Peirce might have written that last paragraph.'' I do not say why. ``More than one path can lead to the Way.''
